% python_api.tex — direct object-level Python API usage (GitTree / GitRepo)

\chapter{Direct Python Object API}
\label{chap:python-api}

This chapter shows how to manipulate \cgspkg{} objects directly (without the
CLI wrapper), including loading a \texttt{.gts} snapshot, reconstructing the
runtime registry, mirroring identity into \texttt{GitTree}/\texttt{GitRepo},
and propagating a tag target.

For normal Multi-repo Sync usage, prefer the CLI lifecycle documented in the
README and Getting Started guide. The reference test topology is CGSil1
(\texttt{https://gitlab.com/CGS\_test/CGSil1}):
\texttt{initialise -> pull -> checkout/add/commit/push -> freeze -> launch\_release}.

The lower-level object/API lifecycle remains:
\texttt{load(.cgs) -> .gts LOADED -> expand(.cgs|.gts) -> .gts PENDING ->
validate(.cgs|.gts) -> .gts READY -> pull(.cgs|.gts) -> checkout(.gts) ->
add -> git(tree,"commit",msg) -> git(tree,"push") -> git(tree,"tag",name) ->
freeze}.

\section{Programmatic Project Configuration}

Project configuration does not require the CLI or an existing \texttt{.cgs}
file. The public facade accepts authoring values non-interactively and delegates
normalization, static validation, and optional serialization to
\texttt{cgs\_format.py}. It performs no Git or network operations.

\begin{lstlisting}[language=Python]
from ComplexGitSync import ComplexGitSyncClient

client = ComplexGitSyncClient()
document = client.configure(
    "GX4G",
    ["codeberg:GX4G/GX4G"],
    output_path="GX4G.cgs",  # optional
)
reference_tree = document.to_git_tree()
\end{lstlisting}

The CLI commands \texttt{configure}, \texttt{create-cgs}, and direct
\texttt{initialise --project/--repo} authoring are adapters over this method.

\section{From Empty Registry to READY via \texttt{.gts}}

\begin{lstlisting}[language=Python]
from pathlib import Path

from ComplexGitSync import GitRepo, GitTree, RefKind, TreeLifecycleState
from ComplexGitSync.orchestre import GtsDocument, build_registry_from_gts_document
from ComplexGitSync.git_tree import WorkingGitTree

# 1) Empty registry lifecycle
empty_registry = WorkingGitTree()
assert empty_registry.recompute_tree_state() == TreeLifecycleState.UNLOADED

# 2) Load the CGSil1 runtime snapshot (.gts) and rebuild the dependency registry
cgshome = Path("../..")
gts_doc = GtsDocument.from_toml(cgshome / ".cgitsync/state/CGSil1.gts")
registry = build_registry_from_gts_document(gts_doc)
assert registry.lifecycle_state == TreeLifecycleState.READY

# 3) Rebuild GitTree/GitRepo identity objects from discovered registry entries
tree = GitTree()
for entry in registry.values():
    tree.add_repo(
        GitRepo(
            project_owner_name=entry.project_owner_name or "owner",
            project_name=entry.project_name or entry.name,
            gitprovider=entry.gitprovider,
            access_protocol=entry.access_protocol,
            commit_sha=entry.commit_sha,
        )
    )
\end{lstlisting}

\section{Object-Level Tag Propagation}

\begin{lstlisting}[language=Python]
# Set the same tag target across all registry entries (parent -> leaves)
tree.propagate_tag(registry, "v1.2.3")
for entry in registry.values():
    assert entry.target_ref_kind == RefKind.TAG
    assert entry.target_ref_name == "v1.2.3"
\end{lstlisting}

This object-level flow is the same state model used by the CLI and
\texttt{ComplexGitSyncClient}, but exposed directly for advanced automation and
tests.

\section{READY Methods in the Client API}

The same lifecycle model is exposed by \texttt{ComplexGitSyncClient}.  In
practice, the workflow is:

\begin{itemize}
  \item \texttt{load(...)} is the step-1 name for direct source loading.
  \item \texttt{expand(...)} is the canonical step-2 name; it runs nested
    \texttt{.cgs} discovery from parents to leaves (recursive).
  \item \texttt{validate(...)} is the canonical step-3 name; it accepts both
    \texttt{.cgs} and \texttt{.gts} sources.
  \item \texttt{initialise(...)} bootstraps a tree and must finish in
    \texttt{READY}; \texttt{clone(...)} is the direct clone entry point for a
    \texttt{.cgs} source; \texttt{pull(...)} resynchronises an existing tree.
  \item \texttt{git(tree, command, *args)} is the unified step-5 interface for
    \texttt{"commit"}, \texttt{"push"}, and \texttt{"tag"} operations.  Each
    follows leaf-first ordering.
  \item \texttt{freeze(...)} emits the next persisted \texttt{.gts} state.
\end{itemize}

\begin{lstlisting}[language=Python]
from pathlib import Path

from ComplexGitSync import ComplexGitSyncClient

client = ComplexGitSyncClient()
cgspath = Path("../..")
cgshome = cgspath / "CGSil1"

# 1. load the CGSil1 test-case .cgs -> .gts LOADED
client.load("../CGSil1.cgs")

# 2. expand(.cgs/.gts) -> .gts PENDING
print(client.expand("../CGSil1.cgs"))

# 3. validate(.cgs/.gts) -> .gts READY
client.validate("../CGSil1.cgs")

# 4. clone(.cgs/.gts)
# Same default as the CLI: output_path defaults to CGSPATH=../..
client.initialise("../CGSil1.cgs")

# Equivalent explicit form:
client.initialise("../CGSil1.cgs", output_path=cgspath)

# Recovery path for partial/stale clone state:
client.clean_initialise_cgs("../CGSil1.cgs", output_path=cgspath)

# Cleanup only:
client.purge_cgs("../CGSil1.cgs", output_path=cgspath)

# 5. READY methods - unified git interface
registry = client.get_dependency_registry()
# Pull uses ROOT -> PARENT -> LEAF; every repository in the tree is a plain
# independent clone (there are no gitlinks) and receives its own git pull.
# A .gts source must be an actual snapshot path (find the current one via
# RuntimeStateStore.latest_snapshot_for(), the same lookup the CLI uses when
# a command's snapshot argument is omitted) -- unlike a .cgs source, there is
# no fixed, predictable filename to hard-code here.
client.pull("../CGSil1.cgs")
# pull() fetches every branch of each remote, so checkout() can join a branch
# somebody else pushed: it starts the local branch from the remote-tracking
# ref when one exists, and only otherwise from the current HEAD.
client.checkout("feature/my-branch")
# Mutations use LEAF -> PARENT -> ROOT.
client.add()
client.git(registry, "commit", "feat: update CGSil1 CGS#1")
client.git(registry, "push")
client.git(registry, "tag", "v1.2.3")

# 6. freeze -> next .gts id
client.freeze("release-2026.05")
\end{lstlisting}

\section{Landing a Branch: \texttt{merge} and \texttt{refresh\_private}}

\texttt{merge(project\_branch, private=False, ff\_only=False, no\_ff=False)}
merges a branch across the tree, leaf-first, and returns one
\texttt{(repo\_name, merged\_ref)} pair per repository a merge moved.

The argument is always the \textbf{project's} branch name. Each repository
resolves what that means for itself, so a private/local configuration
repository merges \texttt{<its default\_branch>\_<project\_branch>} rather
than a branch of the project's name, which it does not have.
\texttt{merge\_plan(project\_branch, private=False)} returns one
\texttt{(repo\_name, source\_ref, status, conflicting\_paths)} row per
repository without merging anything --- what the CLI's \texttt{--dry-run}
prints. \texttt{status} is \texttt{"merge"}, \texttt{"already-on-it"},
\texttt{"no-branch"} or \texttt{"conflicts"}, and the last one carries the
files that block it.

Every repository in scope is checked before any is merged, so a conflict
anywhere leaves the whole tree untouched. The check names the conflicting
files, and counts binary files: a tracked PDF changed on both sides
conflicts just as a source file does.

\texttt{merge\_resolve(project\_branch, private=False, ff\_only=False,
no\_ff=False)} is the opt-in alternative. It merges one repository at a
time, stops at the first that conflicts, and returns a
\texttt{ResolveOutcome} naming what it merged, where it stopped, the files
that stopped it, and the repositories it never reached. It gives up the
all-or-nothing guarantee on purpose: that is the only way to reach the
conflicted worktree a merge tool needs.

\texttt{open\_merge\_tool(repo\_name)} then opens one repository's conflicts
in a merge tool. It returns \texttt{None} once the tool has run, or the
command to run by hand when no tool is available --- a missing editor is a
normal outcome, not an error. A \texttt{merge.tool} you configured always
wins, and your Git configuration is never written.

\texttt{refresh\_private()} is the other direction, and what
\texttt{pull --private} runs: it fetches and merges each private/local
repository's base branch into the derived branch it is on, so a settings
branch does not drift behind the project's while a feature branch is open.

\begin{lstlisting}[language=Python]
client.merge("multi-branch")                  # the project's own repos
client.merge("multi-branch", private=True)    # ComplexGitSync_multi-branch
                                              #   -> ComplexGitSync
client.merge("multi-branch", all_writable=True)   # both halves, one pass
client.refresh_private()                      # the other way round
\end{lstlisting}

\paragraph{\texttt{all\_writable}: the API half of \texttt{--all}.}
\texttt{add}, \texttt{commit}, \texttt{push}, \texttt{merge},
\texttt{merge\_resolve} and \texttt{merge\_plan} each take
\texttt{all\_writable=False} alongside \texttt{private=False}. Setting it
runs the call across this project's own repositories \emph{and} its
writable configuration ones in a single pass. The two are mutually
exclusive and passing both raises; leaving both unset keeps today's
behaviour exactly. Read-only configuration repositories are never written
to under any combination --- the parameter maps to
\texttt{RepoScope.WRITABLE}, which is narrower than
\texttt{RepoScope.ALL}.

\section{Targeting a Single Path: \texttt{add} and \texttt{remove}}

\texttt{add(paths=None)} and \texttt{remove(paths)} accept one or more
filesystem paths (each relative to \texttt{CGSHOME}, or absolute)
instead of operating on the whole tree:

\begin{lstlisting}[language=Python]
# Stage only these two files, each in the repo that owns it -- every
# other repo in the tree is untouched.
client.add(paths=["notes.md", "deps/leaf/src/main.py"])

# Remove a single tracked file: deleted from disk, removal staged.
client.remove(["deps/leaf/obsolete.txt"])
\end{lstlisting}

Both resolve each path to its owning repo through the same helper,
\texttt{resolve\_repo\_for\_path(tree, path)} (\texttt{git\_tree.py}) ---
the deepest (most specific) repo whose \texttt{absolute\_path} contains the
resolved path wins, so a nested child is correctly preferred over its
parent. A path outside every repo in the tree raises
\texttt{GitSyncError} immediately, before anything is staged or removed;
for \texttt{remove}, a path that resolves to a directory or does not exist
raises the same way rather than partially applying. \texttt{add()} with no
\texttt{paths} keeps today's exact behaviour: every repo staged in full,
leaf-first.

\texttt{remove} is a plain \texttt{git rm} (deletes from disk, stages the
removal) --- unrelated to \texttt{GitRunner.rm\_cached} (index-only, used
internally by \texttt{import\_submodules} to drop a gitlink while keeping
the child's working tree); the two are not interchangeable.

\section{\texttt{.gitignore} Lifecycle Sync and Commit Identity}

\texttt{initialise\_cgs}, \texttt{initialise\_cgs\_document},
\texttt{clean\_initialise\_cgs}, \texttt{clean\_init}, \texttt{restart}, and
\texttt{pull} (for \texttt{.cgs} sources) all accept the same four optional
keyword arguments, mirroring the CLI's \texttt{--commit-gitignore},
\texttt{--force-gitignore-sync}, \texttt{--git-user-name}, and
\texttt{--git-user-email} flags:

\begin{lstlisting}[language=Python]
from ComplexGitSync import ComplexGitSyncClient, MasterConfig

client = ComplexGitSyncClient()
client.initialise_cgs(
    "../CGSil1.cgs",
    output_path=cgspath,
    commit_gitignore=True,       # stage/commit/push each changed .gitignore
    force_gitignore_sync=False,  # pull-force fallback for a blocked repo
    git_user_name="cgitsync-bot",
    git_user_email="bot@example.com",
)

# The identity above is persisted to CGSHOME/.cgitsync/master.toml, so a
# later call on the same workspace picks it up without repeating it:
client.initialise_cgs("../CGSil1.cgs", output_path=cgspath, commit_gitignore=True)
assert MasterConfig.resolve_identity(cgshome, client.git_runner) == (
    "cgitsync-bot",
    "bot@example.com",
)
\end{lstlisting}

All four keyword arguments default to \texttt{False}/\texttt{None}: by
default, every repository with children (root, or any nested repository
that itself has further nested children) is safely pulled and has its
\texttt{.gitignore} updated, but nothing is staged, committed, or pushed,
and the commit identity --- were a commit ever made --- would be whatever
\texttt{git config user.name}/\texttt{user.email} already resolves to
locally. \texttt{MasterConfig} (\texttt{master.py}) is the class behind the
identity override: \texttt{configure()} sets it in memory for the current
process, \texttt{persist(cgshome, ...)} writes it to
\texttt{CGSHOME/.cgitsync/master.toml} (and also updates the in-memory
value), \texttt{load(cgshome)} reads a previously persisted override back
in, and \texttt{resolve\_identity(repo\_path, git\_runner)} returns the
\texttt{(user\_name, user\_email)} pair --- or \texttt{(None, None)} to defer
entirely to local git config. This file is workspace-local configuration,
not part of the \texttt{.cgs}/\texttt{.gts} project spec, and is preserved
(not deleted) by \texttt{purge}/\texttt{clean-init}.

\section{Forcing a Clone Protocol}

\texttt{initialise\_cgs}, \texttt{initialise\_cgs\_document},
\texttt{clean\_initialise\_cgs}, \texttt{clean\_init}, \texttt{bootstrap},
and \texttt{clone\_cgs} additionally accept \texttt{force\_access\_protocol},
mirroring the CLI's \texttt{--force-protocol} flag on \texttt{initialise},
\texttt{bootstrap}, and \texttt{clean-init} (\texttt{restart}/\texttt{pull}
do not accept it):

\begin{lstlisting}[language=Python]
# Every repo this run clones -- root siblings and any child discovered
# later from a nested .cgs in a different repo -- is cloned over https,
# regardless of what its own .cgs entry says. No .cgs file is read
# differently or written.
client.initialise_cgs(
    "../CGSil1.cgs",
    output_path=cgspath,
    force_access_protocol="https",
)
\end{lstlisting}

\texttt{"ssh"} or \texttt{"https"}; \texttt{None} (the default) leaves every
entry's own \texttt{.cgs}-declared \texttt{access\_protocol} untouched. The
override lives only in memory for the duration of the call -- it is applied
at the single point every clone's remote URL is built
(\texttt{ComplexGitSyncClient.\_build\_remote\_url}), so it reaches root
siblings and nested-discovered children alike without \texttt{cgs\_format.py}
parsing, or any \texttt{.cgs} file on disk, ever being touched differently.


\section{Repository Discovery: \texttt{discover\_repos}}

\texttt{discover\_repos(root\_dir=None, *, max\_depth=5, output=None)} scans a
directory for git repositories already checked out on disk and drafts a
\texttt{.cgs} from what it finds --- the entry point for adopting a project
that exists but has no specification yet:

\begin{lstlisting}[language=python]
from ComplexGitSync.orchestre import ComplexGitSyncClient

client = ComplexGitSyncClient()

# Read-only scan - nothing is cloned, fetched, or modified
report = client.discover_repos("/path/to/checkout")
for repo in report.repos:
    print(repo.relative_path, repo.identifier, repo.branch)

for warning in report.warnings:
    print("unresolved:", warning)

# Write the draft out (same call, with an output path)
report = client.discover_repos("/path/to/checkout", output="draft.cgs")
\end{lstlisting}

The method returns a \texttt{DiscoverReport} dataclass with fields
\texttt{root} (the scanned directory), \texttt{repos} (a tuple of
\texttt{DiscoveredRepo}), \texttt{cgs\_entries} (authoring-form repo tables
ready for \texttt{configure()}), \texttt{warnings}, and
\texttt{project\_name}.  Each \texttt{DiscoveredRepo} carries
\texttt{relative\_path}, \texttt{absolute\_path}, \texttt{remote\_url},
\texttt{identifier}, \texttt{branch}, and \texttt{has\_cgs}.

Repository addresses are resolved through the existing
\texttt{parse\_repo\_id()} grammar; a repository with no \texttt{origin} or
an unrecognised provider is reported in \texttt{warnings} and omitted from
\texttt{cgs\_entries} rather than guessed at.  Passing \texttt{output}
raises \texttt{GitSyncError} when nothing resolvable was found, so an empty
draft is never written.

\section{Submodule Migration: \texttt{import\_submodules}}

\texttt{import\_submodules(repo\_root, *, apply=False, recursive=False)}
converts git submodules in a local repository to plain ComplexGitSync
nested clones:

\begin{lstlisting}[language=python]
from ComplexGitSync.orchestre import ComplexGitSyncClient

client = ComplexGitSyncClient()

# Dry run - inspect without changing anything
report = client.import_submodules("/path/to/repo")
for sub in report.submodules:
    print(sub.name, sub.path, sub.url)

# Apply - convert every level, root first
report = client.import_submodules(
    "/path/to/repo",
    apply=True,
    recursive=True,
)
print(report.converted)   # tuple of converted submodule names
\end{lstlisting}

The method returns an \texttt{ImportSubmodulesReport} dataclass with fields
\texttt{submodules} (all entries found), \texttt{applied} (whether the
conversion ran), \texttt{converted} (names of converted submodules), and
\texttt{scan\_root} (the repository the call was pointed at, which every
reported path can be expressed from).  When \texttt{apply=False} (the
default), no files are modified.  When \texttt{apply=True}, each submodule
is removed from the index (\texttt{git rm --cached}), its stanza is removed
from \texttt{.gitmodules}, and the parent's \texttt{.gitignore} is updated
to ignore the child path.  The staged changes should then be reviewed and
committed manually.

\texttt{recursive=True} also converts any submodule that has its own
checked-out \texttt{.gitmodules}, at any depth.  That walk follows the
submodule graph declared by \texttt{repo\_root}'s own \texttt{.gitmodules},
so a repository whose holder has already been converted is no longer
reachable --- which is why the conversion of a whole tree happens once, and
last (see \texttt{init\_from\_submodules} below).

\section{Adopting a Submodule Tree: \texttt{init\_from\_submodules}}

\texttt{init\_from\_submodules(repo\_root, *, cgs\_path=None,
max\_depth=5, dry\_run=False, force=False, force\_access\_protocol=None)}
performs the whole adoption of a submodule-based checkout, in the one order
that works:

\begin{lstlisting}[language=python]
from ComplexGitSync.orchestre import ComplexGitSyncClient

client = ComplexGitSyncClient()

# Report the plan without writing, cloning, or converting anything
plan = client.init_from_submodules("/work/project", dry_run=True)

# Adopt: discover -> write .cgs -> initialise -> convert every submodule
report = client.init_from_submodules("/work/project")
print(report.cgs_path)                     # the .cgs used
print(report.tree.is_ready())              # True
print(report.import_report.converted)      # converted submodule names
\end{lstlisting}

Each step delegates to the method that already owns it
(\texttt{discover\_repos}, \texttt{initialise\_cgs},
\texttt{import\_submodules}); the value this method adds is the ordering
and the preconditions.  \texttt{initialise\_cgs} adopts the root in place
but deletes and re-clones every other repository from its remote, where the
submodules are still declared --- so converting first is undone for every
non-root repository, and a repair pass cannot reach any deeper level once
the root's \texttt{.gitmodules} is gone.

The method returns an \texttt{InitFromSubmodulesReport} with fields
\texttt{root} (which is also \texttt{CGSHOME}), \texttt{discover} (the
\texttt{DiscoverReport} from step~1), \texttt{cgs\_path},
\texttt{cgs\_written}, \texttt{import\_report}, \texttt{tree}, and
\texttt{dry\_run}.  Nothing is committed: the conversion touches every
repository that held a submodule, so \texttt{branch}, \texttt{checkout},
\texttt{add}, and \texttt{commit} stay explicit, separate calls.  A
\texttt{GitSyncError} is raised before anything is written when the
discovery resolved nothing, when the project name does not match the
directory name (\texttt{CGSHOME} is \texttt{<parent>/<project-name>}), or
when \texttt{repo\_root} has no \texttt{.gitmodules} to convert
(\texttt{force=True} overrides the last one).

\section{Looking at a memory: \texttt{memory\_status},
\texttt{memory\_list}, \texttt{memory\_show}}

The three read-only questions the CLI's \texttt{memory} group asks, each
returning plain data rather than text:

\begin{lstlisting}[language=python]
from ComplexGitSync.orchestre import ComplexGitSyncClient

client = ComplexGitSyncClient()

status = client.memory_status("/path/to/CGSHOME")
status["states"], status["entries"]      # how much is remembered
status["verification"]                   # verified / no-history / legacy / corrupt
status["latest_toolchain"]["git"]        # what produced the newest record

for row in client.memory_list("/path/to/CGSHOME"):
    row["state"], row["recorded_at"], row["commands"]

state = client.memory_show("/path/to/CGSHOME", "2acdc98b")   # hash or prefix
state["repos"], state["entries"]
\end{lstlisting}

\texttt{memory\_show} raises \texttt{GitSyncError} when the prefix matches
no State, or more than one. \texttt{memory\_list} includes a State that
entries name but that is not on disk, with \texttt{path} set to
\texttt{None} --- \texttt{verify} reports it as \texttt{MISSING\_STATE},
and this is how a caller finds out which one.

\section{Machine-readable output: \texttt{status\_json} and
\texttt{verify\_json}}

Two methods return the same answers as \texttt{status()} and
\texttt{verify()}, serialised as one JSON object each. They are what the
CLI's \texttt{--json} flag prints, so a Python caller and a shell pipeline
read exactly the same fields.

\begin{lstlisting}[language=python]
import json
from ComplexGitSync.orchestre import ComplexGitSyncClient

client = ComplexGitSyncClient()
client.load_gts("/path/to/snapshot.gts")

payload = json.loads(client.status_json())
payload["cgitsync_branch"]          # the branch the tree is on
payload["summary"]["repos"]         # how many repositories it holds
payload["repositories"][0]["sync"]  # per repository, as the table shows it

chain = json.loads(client.verify_json("/path/to/CGSHOME"))
chain["status"]                     # "clean" or "findings"
\end{lstlisting}

\texttt{verify\_json} keeps the report it built on
\texttt{client.last\_verify\_report}, so a caller needing both the object
and the verdict does not verify the chain twice --- which under
\texttt{repair=True} would mean repairing twice.

The shape of both objects lives in \texttt{json\_render.py}, defined once
for every command rather than per command, and carries a
\texttt{schema\_version}. The promise attached to it is \textbf{additive
only}: new fields may appear; existing fields do not change meaning and do
not vanish.

\section{A memory as a repository: \texttt{memory\_init},
\texttt{memory\_clone}, \texttt{memory\_push}}

\begin{lstlisting}[language=python]
from ComplexGitSync.orchestre import ComplexGitSyncClient

client = ComplexGitSyncClient()
client.load_gts("/path/to/CGSHOME/.cgitsync/state/<hash>.gts")

proposal = client.memory_init("/path/to/CGSHOME")
proposal["entry"]        # the .cgs entry that mounts the memory
proposal["branch"]       # this project's branch of the memory repository
proposal["create_with"]  # the command that creates it -- never run for you

client.memory_push("/path/to/CGSHOME")                      # commit + push
client.memory_clone("/other/workspace", owner="flipoyo")    # onto a new machine
\end{lstlisting}

\texttt{memory\_clone} needs no loaded project --- that is the case it
exists for --- so give it an \texttt{owner} and a \texttt{branch} when
nothing is loaded, or a \texttt{remote} address directly for a memory that
does not live where the convention says. It raises
\texttt{GitSyncError} rather than writing over a memory already present, or
when the branch it wants has never been pushed.

\section{Register Integrity: \texttt{verify}}

\texttt{verify(cgshome, *, repair=False)} checks the hash-chained
\texttt{.cgitsync/lgr} register under \texttt{cgshome} for tamper-evidence:

\begin{lstlisting}[language=python]
from ComplexGitSync.orchestre import ComplexGitSyncClient

client = ComplexGitSyncClient()
report = client.verify("/path/to/CGSHOME")

if report.is_clean:
    print("register is clean")
else:
    for seq, finding, detail in report.findings:
        print(seq, finding.name, detail)

# Repair a stale HEAD cache (never rewrites or deletes an entry)
client.verify("/path/to/CGSHOME", repair=True)
\end{lstlisting}

Returns an \texttt{integrity.VerificationReport}. Read
\texttt{report.state} first: it is one of four \texttt{HistoryState}
members --- \texttt{VERIFIED}, \texttt{NO\_HISTORY}, \texttt{LEGACY},
\texttt{CORRUPT} --- and \texttt{report.is\_verified} is \texttt{True}
only for the first.

\texttt{report.is\_clean} is a narrower question: it means ``no findings
were recorded'', which an empty register also satisfies. Use it to ask
whether anything is \emph{wrong}; use \texttt{state} to ask what was
actually established. \texttt{findings} holds \texttt{(seq, Finding,
detail)} tuples.

Every command that writes a State now appends an entry to
\texttt{.cgitsync/lgr/}, so a workspace used since that landed answers
\texttt{VERIFIED}. A workspace whose only history predates it answers
\texttt{LEGACY}: the single-file register is still read, and nothing writes
it any more.  \texttt{Finding} enumerates
\texttt{BROKEN\_LINK}, \texttt{BAD\_ENTRY\_HASH}, \texttt{SEQ\_GAP},
\texttt{SEQ\_DUPLICATE}, and \texttt{HEAD\_STALE} for what this method
actually checks (chain linkage and the \texttt{HEAD} cache); its remaining
three members (\texttt{MISSING\_STATE}, \texttt{ORPHAN\_STATE},
\texttt{STATE\_DIGEST\_MISMATCH}) are reserved for a future store-level
verification pass that cross-references entries against the actual
\texttt{state/<hash>.gts} files on disk.

With \texttt{repair=True}, a stale \texttt{HEAD} cache is corrected in
place. Register entries are never rewritten or deleted by this method,
repair or not — a broken chain is reported, not silently healed.
